
\medskip

Cet exercice est un questionnaire à choix multiple. Pour chaque question, une seule des trois propositions est exacte.

Indiquer sur la copie le numéro de la question et la lettre de la proposition
choisie.

Aucune justification n'est demandée.

\medskip

\emph{Pour chaque question, une réponse exacte rapporte un point.\\
Une réponse fausse, une réponse multiple ou l'absence de réponse ne rapporte ni n'enlève de point.}

\medskip

Dans toutes les questions suivantes, l'espace est rapporté à un repère orthonormé.

\medskip

\begin{enumerate}
\item %On considère la droite $\Delta_1$ de représentation paramétrique 
%$\left\{\begin{array}{l c l}
%x &=& 1 - 3t\\
%y &=& 4 + 2t\\
%z &=&\phantom{4 +2}t
%\end{array}\right.$, où $t \in  \R$
%ainsi que la droite $\Delta_2$ de représentation paramétrique 
%$\left\{\begin{array}{l c l}
%x& =& -4 + \phantom{2}s\\
%y& =& \phantom{-}2 + 2s\\
%z& =& - 1 + \phantom{2}s
%\end{array}\right.$, où $s \in \R$.
%\begin{enumerate}
%\item Les droites $\Delta_1$ et $\Delta_2$ sont parallèles.
%\item Les droites $\Delta_1$ et $\Delta_2$ sont orthogonales.
%\item Les droites $\Delta_1$ et $\Delta_2$ sont sécantes.
%\end{enumerate}
$\bullet~~$ $\Delta_1$ a pour vecteur directeur $\delta_1\begin{pmatrix}-3\\2\\1\end{pmatrix}$ et $\Delta_2$ a pour vecteur directeur $\delta_2\begin{pmatrix}1\\2\\1\end{pmatrix}$ : les premières coordonnées ne sont pas égales, les vecteurs directeurs ne sont pas colinéaires, les droites $\Delta_1$ et $\Delta_2$ ne sont pas parallèles ; 

%$\bullet~~$ On a $\delta_1 \cdot \delta_2 = -3 \times 1 + 2 \times 2 + 1 \times 1 = 2 \ne 0$ : les vecteurs directeurs ne sont pas orthogonaux, les droites ne sont pas perpendiculaires.
%
%$\bullet~~$Si les droites sont sécantes les coordonnées de leur point commun vérifient le système :
%
%$\left\{\begin{array}{l c l}
%1 - 3t	&=&-4 + s\\
%4 + 2t	&=&2 + 2s\\
%t		&=&- 1 + s
%\end{array}\right.$

%En remplaçant $t$ par $- 1 + s$ dans les deux premières équations on obtient :
%
%$\left\{\begin{array}{l c l}
%1 + 3 - 3s&=&-4 + s\\
%4  - 2 + 2s&=&2 + 2s\\
%t&=&- 1 + s
%\end{array}\right.$, soit 
%$\left\{\begin{array}{l c l}
%1 + 3 - 3s&=&-4 + s\\
%4 + - 2 + 2s&=&2 + 2s\\
%t&=&- 1 + 2s
%\end{array}\right. \iff 
%\left\{\begin{array}{l c l}
%8&=&4s\\
%2 + 2s &=&2 + 2s\\
%t&=&- 1 + 2s
%\end{array}\right. \iff \left\{\begin{array}{l c l}
%2&=&s\\
%2 + 2s &=& 2 + 2s\\
%t&=&- 1 + 2s
%\end{array}\right.\iff \left\{\begin{array}{l c l}
%2&=&s\\
%2 + 2s &=& 2 + 2s\\
%t&=&-1
%\end{array}\right.$

En remplaçant respectivement dans les équations paramétriques $s$ par 2 et $t$ par 1, on trouve que le point de coordonnées $(-2~;~6~;~1)$ appartient bien aux deux droites.
\item %On considère la droite $d$ de représentation paramétrique 
%$\left\{\begin{array}{l c l}x&=&1 + \phantom{2}t\\
%y &=& 3 - \phantom{2}t\\
%z &=& 1 + 2t
%\end{array}\right.$, où $t \in R$,
%et le plan $P$ d'équation cartésienne : $4x + 2y - z + 3 = 0.$
	\begin{enumerate}
		\item %La droite $d$ est incluse dans le plan $P$.
$M(x = 1 + 2t~;~y = 3 - t~;~z = 1 + 2t) \in d$ appartient à $P$ si ses coordonnées vérifient l'équation cartésienne de $P$, donc si 

$4(1 + t) + 2(3 - t) -(1 + 2t) + 3 = 0 \iff 4 + 4t + 6 - 2t - 1 - 2t + 3 = 0 \iff 0t + 12 = 0 \iff 0t = -12$ : cette équation n'a pas de solution : cela signifie que la droite $d$ est parallèle strictement au plan $P$.

%		\item %La droite $d$ est parallèle strictement au plan $P$.

%		\item %La droite $d$ est sécante au plan $P$.
	\end{enumerate}
\item On considère les points A(3~;~2~;~1), B(7~;~3 ~;~1), C$(-1~;~4~;~5)$ et D$(- 3~;~3~;~5)$.

Commençons par la deuxième affirmation : on a 

$\vect{\text{AB}}\begin{pmatrix}4\\1\\0\end{pmatrix}, \: \vect{\text{AC}}\begin{pmatrix}-4\\2\\4\end{pmatrix}$ : ces vecteurs ne sont pas colinéaires, donc A, B et C ne sont pas alignés.

$\vect{\text{CD}}\begin{pmatrix}-2\\-1\\0\end{pmatrix}$ : ces vecteurs ne sont pas colinéaires.

Il reste donc la première affirmation que l'on peut démontrer :

A, B, C et D sont coplanaires, c'est-à-dire si par exemple D appartient au plan (ABC) et dans ce cas si $\vect{\text{AD}} = x\vect{\text{AB}} + y\vect{\text{AC}}$, avec $x \in \R,\: y \in \R$ (car $\vect{\text{AB}}$ et $\vect{\text{AC}}$ ne sont pas des vecteurs colinéaires.

Avec $\vect{\text{AD}}\begin{pmatrix}-6\\1\\4\end{pmatrix}$, l'égalité vectorielle précédente se traduit par le système :

$\left\{\begin{array}{l c r}
-6&=&4x - 4y\\
1&=&x + 2y\\
4 &=& 4y
\end{array}\right. \iff \left\{\begin{array}{l c r}
-6&=&4x - 4\\
1&=&x + 2\\
1 &=& y
\end{array}\right. \iff
\left\{\begin{array}{l c r}
-2&=&4x\\
-1&=&x\\
1 &=& y
\end{array}\right. \iff \left\{\begin{array}{l c r}
-\frac12&=&x\\
-1&=&x\\
1 &=& y
\end{array}\right.$

Ce système n'a pas de solution. Donc le point D n'appartient pas au plan défini par les points A, B et C ou encore A, B, C et D ne sont pas coplanaires.
%	\end{enumerate}
\item %On considère les plans $Q$ et $Q'$ d'équation cartésienne respective $3x - 2y + z + 1 = 0$ et $4x + y - z + 3 = 0$.
	\begin{enumerate}
		\item %Le point R$(1~;~1~;~-2)$ appartient aux deux plans.
R$(1~;~1~;~-2) \in (Q) \iff 3 \times 1 + -2  \times1  - 2 +1 = 0 \iff 3 - 2 - 2 + 1 = 0$ qui est vraie ;
		
R$(1~;~1~;~-2) \in (Q') \iff 4 \times 1 + 1 \times 1  - (- 2) + 3 = 0 \iff 10 = 0$ qui est fausse : la première affirmation est fausse ;
		\item %Les deux plans sont orthogonaux.
$(Q)$ a pour vecteur normal $\vect{q}\begin{pmatrix}3\\- 2\\1\end{pmatrix}$ ; 
$(Q')$ a pour vecteur normal $\vect{q'}\begin{pmatrix}4\\1\\- 1\end{pmatrix}$ ; 

$\vect{q} \cdot \vect{q'} = 3 \times  4 -2 \times  1 + 1 \times (-1) = 12 - 2 - 1 = 9 \ne 0$ : les vecteurs normaux ne sont pas orthogonaux, les plans $(Q)$ et $(Q')$ ne sont pas perpendiculaires.
		\item %Les deux plans sont sécants avec pour intersection la droite de représentation
%paramétrique $\left\{\begin{array}{l c l}
%x&=& \phantom{17}t\\
%y& =& \phantom{1}7t + 4 \\
%z& =& 11t + 7
%\end{array}\right.$, où $t \in \R$.

Les points communs aux deux plans sont les points dont les coordonnées vérifient les équations de ces deux plans donc le système :

$\left\{\begin{array}{l cl}
3x - 2y + z + 1 &=& 0\\
4x + y - z + 3 &=& 0
\end{array}\right.$

On peut poser $x = t$, avec $t \in \R$ et on obtient le système aux deux inconnues $y$ et $z$ suivant :

$\left\{\begin{array}{l cl}
x&=&t\\
3x - 2y + z + 1 &=& 0\\
4x + y - z + 3 &=& 0
\end{array}\right. \iff \left\{\begin{array}{l cl}
x&=&t\\
-2y + z &=& - 3t - 1\\
\phantom{- 2}y - z &=& - 4t - 3
\end{array}\right. $

La somme membre à membre des deux dernières équations donne

$-y =- 7t - 4 = 0 \iff y = 7t + 4 $ ;

enfin la deuxième équation $-2y + z = - 3t - 1$ donne $z = 2y - 3t - 1 = 2(7t + 4) - 3t - 1 = 14t + 8 - 3t - 1 = 11t + 7$.
		\end{enumerate}


Conclusion : les points communs aux deux plans sont les points dont les coordonnées vérifient le système : $\left\{\begin{array}{l c l}
x&=& \phantom{17}t\\
y& =& \phantom{1}7t + 4 \\
z& =& 11t + 7
\end{array}\right.$, où $t \in \R$ : c'est l'équation de la droite contenant le point K(0~;~4~;~7) et dont un vecteur directeur est le vecteur $\vect{u}\begin{pmatrix} 1\\7\\11\end{pmatrix}$. : la troisième affirmation est vraie.
\end{enumerate}

\vspace{0,5cm}

